%!TEX root = main.tex

Let $M$ be the Firrtl program before applying the ExpandConnections pass, and \expandconnects $(M)$ be the program returned.

\paragraph*{Specification}
The correctness of \expandconnects is specified by the formula
%
$$(tm,\,cm) \models \, \expandconnects(M) \rightarrow (tm,\,cm) \models \,M. $$
%
The formula states that if the pair of the typing mapping $tm$ and the
connection mapping $ct$ conform to the program of produce by \expandconnects $(M)$, then the same pair of $tm$ and $ct$ should
also conform to the original program $M$.  More specifically, the
Firrtl programs before and after the ExpandConnections pass should
represent the same low-level circuit.  The process of
ExpandConnections expects an equivalent representation of the
original program, with the aggregate-typed objects now represented in
ground types.

To prove that conformance holds for the original Firrtl module, we need
to derive clues from the conformance after the expansion. This
information can be drawn from the \expandconnects \,function
by cases analysis of different kinds of statements in Firrtl modules.
Consider the connection statement between two references
\mbox{\prettyll{l <= r}},
we illustrate the proof structure for finding the
conformance between the program and the $tm$ and $cm$ representations.

Since connections have no effect on changing the typing mapping
(i.e., $tm$ stays the same), we can focus on the semantic rules stated for $cm$ (as shown in
Fig \ref{fig:firrtl-sem}).
Recall the semantic rule of ``connection:ref'',
which focuses on three predicates.
The first predicate ensures that the orientations of connected elements
fit together, such as from \source\, to \sink\.%($P_{orient}$).
The second predicate ensures the type compatibility of
the connected objects.% ($P_{compat}$).
The last predicate ensures that the connection
mapping $cm$ contains the corresponding connection expressions.% ($P_{findCM}$).

Note that semantics reflect on $cm$, which contains only ground type
elements. This means the three predicates should correspond to a set of
predicates on ground type elements.  They can be named
$P_{orient}$, $P_{compat}$ and $P_{findCM}$, which are the predicates
on two ground type elements. For example, $P_{orient}$ states the
orientations of two atomic elements fit together. %%  Additionally, the
%% connections have no effect on the typing mapping, therefore $tm$ stays
%% the same.

The semantic rule of ``connection:ref'' can represent the conformance
for an aggregate type connection. The atomic elements involved in the connection
should satisfy the predicates modeled with $cm$.
By contrast, the result
of \expandconnects\, is a sequence of connection statements on ground
types. If we directly use the semantic rule for statement sequences, we
need to reason about every intermediate state $cm_n$ providing by the whole
sequence.
%% , like
%% $$
%% cm_0 \xrightarrow[tm, tmap]{\expandconnects(M)} cm_n
%% = cm_0 \xrightarrow[tm, tmap]{s_0} cm_1 \xrightarrow[tm, tmap]{s_1} \dots cm_n.
%% $$
It is more convenient to have one conclusion for the entire connection sequence produced by \expandconnects.

Concretely, we have the following lemmas.
First of all, we prove :
\begin{lemma}
 The result of \expandconnects $(r<=r')$ contains only ground type connections.
\end{lemma}

Then, we state that:
\begin{lemma}\label{p1-p3}
For all connection statements produced by \expandconnects $(r<=r')$, $rl$ and $rl'$ are the lists of ground type references obtained by \expandref$(r)$ and \expandref(r') respectively.\\
Then $all2(P_{orient}, rl, rl')\;\wedge \;all2(P_{compat}, rl,rl')\; \wedge \;all2(P_{findCM}, rl, rl')$.
\end{lemma}
%%
Here, {\sf all2} function is used to check
whether the predicate is satisfied by every pairwise ground
type elements in $rl$ and $rl'$.
%The predicates are the ones given for a pair ground type elements between a connection.
Since the definition of \expandconnects\, for
connections contains checking orientations, the first part $all2(P_{orient}, rl, rl')$ can be  proven.
The type compatibility predicate $all2(P_{compat}, rl,rl')$ can be proved by the type equivalent checking and the flipped fields checking in \expandconnects.
The last $all2(P_{findCM}, rl, rl')$ states that in the final $cm$, we can find the reference in $rl'$ corresponding to the one in $rl'$. This can be proven by the lemma provide by maps data structure, which ensures that the pairs of update and find actions satisfy this condition.
%% Meanwhile, the $all2$ function also satisfies the the size of the two lists are equal.

We also need to prove that the \expandref\, and \listgexp\, produce the same lists for both $r$ and $r'$.
\begin{lemma}\label{rl-rl'}
For all references $r$, the reference expressions obtained by \expandref(r) contain the same reference identifiers generated by \listgexp$(r)$.
\end{lemma}
This can be proved by comparing the definitions of the two functions. They both rely on the type of the reference input.
At the end, we can conclude that the result of \expandconnects $(r<=r')$ implies the three predicates hold as lemma\ref{p1-p3}. With lemma \ref{rl-rl'} we establish that all the predicates satisfied by the lists produced by \expandref\, can also be satisfied by the lists produced by \listgexp. Therefore, the lists produced by \listgexp and the predicates they satisfy exactly match the requirements of the semantic rule of ``connection:ref''
for the original statement $(r<=r')$.

Similarly, we can changing from semantic reasoning on sequences of
statements to applying the predicates on the lists of ground type
elements expanded from each kind of statement. This transition ensures
that we analyze and verify the predicates directly on the lists of
ground type elements generated by functions like \expandref\,
and \listgexp. This approach aligns with proving the conformance of
connections and other statements in the context of Firrtl modules.


%% We can benefit from using the lemmas provided by the Seq library from mathcomp \xiaomu{references}.


\hide{
Let $ss$ be the statement sequence before applying the function {\tt expand\_cnct} and $ss_{\rm cnct}$ be the sequence of connect statements in $ss$.

\paragraph*{Specification}
Then the correctness of {\tt expand\_cnct} is specified by the formula
%
$$\mbox{\tt type\_equivs}(\mbox{\tt expand\_cnct}(ss_{\rm cnct})) \wedge \mbox{\tt all\_ground}(\mbox{\tt expand\_cnct}(ss_{\rm cnct})).$$
%
The formula $\mbox{\tt type\_equivs}(\mbox{\tt expand\_cnct}(ss_{\rm cnct}))$ specifies that for each connect statement generated by $\mbox{\tt expand\_cnct}(ss_{\rm cnct})$, the types of its left-hand-side and right-hand-side expressions are equivalent. On the other hand, the formula $\mbox{\tt all\_ground}(\mbox{\tt expand\_cnct}(ss_{\rm cnct}))$ specifies that all the expressions occurring in the connect statements generated by $\mbox{\tt expand\_cnct}(ss_{\rm cnct})$ have ground types.

More specifically,
the predicate $\mbox{\tt type\_equivs}(ss')$ is recursively defined and requires that each connect statement  $x <= e$ in  $ss'$ satisfies that $\mbox{\tt type\_equiv}(\mbox{\tt type\_of\_expr}(x), \mbox{\tt type\_of\_expr}(e))$ hods, where ${\tt type\_equiv}(t_1, t_2)$ is the predicate that specifies the type equivalence of $t_1$ and $t_2$. On the other hand, the predicate $\mbox{\tt all\_ground}(ss')$ is recursively defined and requires that each connect statement $x <= e$ in  $ss'$ satisfies that $\mbox{\tt ground}(\mbox{\tt type\_of\_expr}(x)) \wedge \mbox{\tt ground}(\mbox{\tt type\_of\_expr}(e))$ holds, where the predicate $\mbox{\tt ground}(t)$ checks that $t$ is a ground type.

\paragraph*{Verification}
Recall that {\tt expand\_cnct} expands a connect statement between components of aggregate types by using the function {\tt expand\_expr} to expand each of the left-hand-side and right-hand-side expressions of the statement into a sequence of expressions of ground types. Since $\mbox{\tt expand\_expr}(e)$ expands the expression $e$ by following the structure of $\mbox{\tt type\_of\_expr}(e)$, we define a recursive function $\mbox{\tt expand\_type}(t)$ that expands an aggregate type $t$ into a list of ground types. Then we prove the following two lemmas to prove $ce \vdash \mbox{\tt type\_equivs}(\mbox{\tt expand\_cnct}(ss_{\rm cnct}))$.

\begin{lemma}\label{expand-expr-expand-type}
For every connect statement $x <= e$ in $ss_{\rm cnct}$,
$$
\begin{array}{l c l}
ce & \vdash & \mbox{\tt type\_equiv}(\mbox{\tt type\_of\_expr}(x), \mbox{\tt type\_of\_expr}(e)) \rightarrow \\
& & \ \ {\tt all2}(\mbox{\tt type\_equiv},{\tt expand\_type}({\tt type\_of\_expr}(x)),{\tt expand\_type}({\tt type\_of\_expr}(e))).
\end{array}
$$
\end{lemma}
In Lemma \ref{expand-expr-expand-type}, the ${\tt all2}$ function is used to check whether the predicate {\tt type\_equiv} is satisfied by every pairwise types in ${\tt expand\_type}({\tt type\_of\_expr}(x))$ and ${\tt expand\_type}({\tt type\_of\_expr}(e))$.

\begin{lemma}\label{expand-expr-expand-expr-type}
For every expression $e$,
$$ce \vdash {\tt map}({\tt type\_of\_expr}, {\tt expand\_expr}(e)) = {\tt expand\_type}\,({\tt type\_of\_expr}(e)).$$
\end{lemma}
Note that in Lemma~\ref{expand-expr-expand-expr-type}, the ${\tt map}$ function is used to apply the function ${\tt type\_of\_expr}$ to every expression in ${\tt expand\_expr}(e)$ to produce a list of types.

Finally, $ce \vdash \mbox{\tt all\_ground}(\mbox{\tt expand\_cnct}(ss_{\rm cnct}))$ can be proved by the following lemma.
\begin{lemma}\label{expand-expr-ground}
For every expression $e$,
$$ce \vdash {\tt all}({\tt ground}, {\tt expand\_type}({\tt type\_of\_expr}(e))).$$
\end{lemma}
Note that in Lemma~\ref{expand-expr-ground}, the ${\tt all}$ function is used to check whether the predicate {\tt ground} is satisfied by every type in ${\tt expand\_type}({\tt type\_of\_expr}(e))$.
}

%%%%%% original texts by xiaomu
%%%%%% original texts by xiaomu
\hide{
Recall the the ExpandConnects pass is aimed to decompose the connect statements of aggregate
types to ground type connections.
To prove the functional correctness of the ExpandConnects pass, we need to ensure
that the sub elements of the components are connected to their equivalent ground types.
Let $ce^{in}$ be the input component environment,
%% such that $ce^{in}$ is the result of the InferTypes pass,
and $ss^{in}$ be the list of connect statements in the input circuit
diagram $D_M^{in}$.
%% Let $ss^{out}$ be the output list of connect
%% statements after applying the Expand Connections pass.
The input
$D_M^{in}$ is required to be the result of the InferTypes pass, in order
to make sure that the connections are made between equivalent types.

We specify the correctness based on the types of the left-hand side and the right-hand side
expressions in the resulting connect statements.
First, we state that the type equivalence holds for the output statements generated by the ExpandConnects pass
as,
\begin{lemma}\label{expand-conct-type1-ss}
$\forall ce^{in}\; ss^{in},\, ce^{in} \models TypeEquivs\, ss^{in} \Rightarrow TypeEquivs\, [{\tt expand\_cnct}\,s\,|\,s\leftarrow ss^{in}].$
\end{lemma}
Intuitively, it states that if the type equivalence predicate holds for all the
connect statements in $ss^{in}$,
it holds for the statements produced by applying the {\tt expand\_cnct} \eqref{expand-cnct} function on all the inputs.
More specifically, the type equivalence predicate $TypeEquivs$ holds for a list of statements if
for each individual comparison of the types on the two sides of a connection, the type equivalence
predicate $TypeEquiv$ holds.
If there are $n$ statements $s_1,s_2,\dots ,s_n$ in the sequence $ss$ of statements,
then $TypeEquivs\;ss\,\!=\,\!\bigwedge_{i=0}^{n} TypeEquiv\;s_i$.

The proof can be achieved by induction on the list of statements,
assuming we have the proof for expanding one statement.
Since the function {\tt expand\_cnct} only affect connect statements,
we can state the following lemma:
\begin{lemma}\label{expand-conct-type1-s}
$\forall ce^{in}\;v\;e,\, ce^{in} \models TypeEquiv\,(v<=e) \Rightarrow TypeEquivs\, ({\tt expand\_cnct}\,(v<=e)).$
\end{lemma}
If the predicate $TypeEquiv$ holds for a connect statement $v<=e$ under a component environment,
then the types of $v$ and $e$ are equivalent.
Lemma \ref{expand-conct-type1-s} can be further simplified as,
\begin{lemma}\label{expand-conct-type1-s}
$\forall ce^{in}\;v\;e,\, ce^{in} \models Equiv\,({\tt typeof}\,v,\;{\tt typeof}\,e) \Rightarrow$
$Equivs\,[({\tt typeof}\,l,\,{\tt typeof}\,r)\,|\,l \leftarrow{\tt expand\_expr}\,v,\;r\leftarrow {\tt expand\_expr}\,e]$.
\end{lemma}
If there are $n$ pairs of types $(t_1^l,t_1^r),(t_2^l,t_2^r),\dots ,(t_n^l,t_n^r)$ in the list $ts$,
then $Equivs$ holds for $ts$ if
$\bigwedge_{i=0}^{n} Equiv\;(t_i^l\,t_i^r)$.
In Lemma \ref{expand-conct-type1-s}, the list of pairs are the types obtained by
expanding the reference on the left-hand side
and expanding the expression on the right-hand side, respectively.
If we can build the relation between the aggregate type of an expression and the types of
the corresponding expanded expressions, the above lemma can be proved.
Here we introduce a function {\tt dec\_type} that can decompose the aggregate types
by its structure to get a list of ground types.
A related lemma is,
\begin{lemma}\label{dec-type}
$\forall ce^{in}\, t_1\, t_2,\,ce^{in}\models Equiv\,(t_1,\;t_2) \Rightarrow$
$Equivs\,[(x,\,y)\,|\,x \leftarrow{\tt dec\_type}\,t_1,\;y\leftarrow {\tt dec\_type}\,$ $t_2]$
\end{lemma}
It states that, decomposing a pair of equivalent aggregate types produces two pairwise equivalent type lists.
If we use the {\tt dec\_type} function on the aggregate type of the expression  to be expanded, it
returns the ground types that are identical to the ones produced by {\tt expand\_expr}, as in the following lemma.
\begin{lemma}\label{expand-expr-dec-type}
$\forall ce^{in}\, e,\,ce^{in}\models[{\tt typeof}\,x\,|\,x\leftarrow{\tt expand\_expr}\,e] = {\tt dec\_type}\,({\tt typeof}\,e)$
\end{lemma}
By rewriting Lemma \ref{expand-expr-dec-type} in the goal of Lemma \ref{expand-conct-type1-s}, the proof goal is change to
$Equivs\,[(x,\,y)\,|\,x\leftarrow{\tt dec\_type}\;({\tt typeof}\,v),\;y\leftarrow {\tt dec\_type}\,({\tt typeof}\,e)]$.
Now we can apply Lemma \ref{dec-type} with assigning $t_1$ and $t_2$ with ${\tt typeof}\,v$ and ${\tt typeof}\,e$
to complete the proof.

Secondly, we state that all the expressions after applying the ExpandConnects pass
have ground types.
\begin{lemma}
$\forall ce ^{in}\,ss^{in},\,ce^{in}\models AllGround ({\tt expand\_cnct}\,ss^{in})$
\end{lemma}

The above lemma states the predicate $AllGround$ holds for the generated
statements of {\tt expand\_cnct}.
Similarly, the proof goal will be reduced to prove that the $AllGround$ predicate holds for the expanded expressions.
It is rather easy to use Lemma \ref{expand-expr-dec-type} we had for {\tt dec\_type}.
\fu{To relate the types of the expanded expressions are identical to decompose the type of the original expression.}
By induction on the aggregate type constructors, it is able to prove that decomposed types by
the {\tt dec\_type} function will always return ground types.
}
%%%%%% original texts by xiaomu
%%%%%% original texts by xiaomu
